\begin{proposition}[The topology of the two modified pinchings]\hypertarget{tag-02ZR}{}\label{push:src:A.4.1}\label{push:A41}
\StatementRecord{push:src:A.4.1}{proved}{02ZR}{\SeriesRef{fam:src:E.1.1}, \SeriesRef{fam:src:E.1.2}, \SeriesRef{fam:src:E.1.3}, \SeriesRef{fam:src:E.1.4}, \SeriesRef{push:src:A.2.2}}

With integral coefficients, the spaces constructed above have
\[
\begin{array}{c|rrrrrrr|r}
&b_0&b_1&b_2&b_3&b_4&b_5&b_6&\chi\\\hline
R&1&0&11&0&11&0&1&24\\
R^\pm&1&0&12&0&12&0&1&26\\
Q&1&1&8&7&10&0&1&12\\
Q^\pm&1&1&9&7&11&0&1&14
\end{array}
\]
All these homology groups are torsion free. The fundamental
groups of $R,R^\pm$ are trivial and those of $Q,Q^\pm$ are
$\mathbb Z$. The dual complex of either $Q^\pm$ is a circle.
\end{proposition}


\begin{proof}
\begin{enumerate}[label=\arabic*.,leftmargin=*]
\item Proposition~\ref{fam:src:E.1.4} gives the integral groups and fundamental groups of $R,Q$, together with the normalization maps $d_i=s_{0*}-s_{\infty*}T_*$. We use exactly its lattice $\mathcal L=H_2(S;\mathbb Z)$, $F=[S_t]$, $m=P-O$, and set $v=v_0=[6P]-[O]-3F$. In these coordinates its map $a\mapsto(aF,av,am)$ is onto $\mathbb Z^3$, its $d_2$ has primitive image of rank three, and its $d_4$ is primitive and injective.

\item Now let $h:R^+\to R$ be the modification constructed in \ref{push:src:A.2.2}.
Its fibres are points off the nodal curve $C=S_\infty$, one
$\mathbb P^1$ over $C_{\mathrm{reg}}$, and a chain of two
$\mathbb P^1$'s over the node. The chart calculation identifies
the sheaf $R^2h_*\mathbb Z$ with
$\eta_*\mathbb Z_{\mathbb P^1}$, where
$\eta:\mathbb P^1\to C$ is the normalization. If $(a,b)$ are the cohomological
degrees on the two exceptional curves, generization to the
two branches is $a$ and $a+b$: on one branch the generic
exceptional fibre specializes to the first curve, on the
other it specializes to their sum. The matrix
$\left(\begin{smallmatrix}1&0\\1&1\end{smallmatrix}\right)$
is unimodular. These generizations follow by substituting
$X=za$, $v=aB$ in the blow-up charts, or by taking the two
limits of their exceptional fibres. Proper local neighbourhoods
retract onto these fibres by the positive toric contraction,
so the fibre calculation computes the higher direct images.
There is no other positive-degree direct image.

\item The Leray spectral sequence has only the row $H^*(R,\mathbb Z)$
and the row $H^{*-2}(\mathbb P^1,\mathbb Z)$; both are free
and in even degrees. It degenerates integrally for degree
reasons, giving the stated groups of $R^+$. The same local
neighbourhood retractions and Van Kampen show that $h$ induces
an isomorphism on fundamental groups, so $R^+$ is simply
connected. The case $R^-$ interchanges the branches.

\item For $Q^+$ use the normalization mapping cylinder once again.
The new degree-two difference map is the old $d_2$ with
one extra coordinate $aF$: the exceptional divisor restricts
to $F$ on the blown-up section and to zero on the other.
It has the same primitive rank-three image and the same
kernel; its cokernel rank increases by one. The degree-four
difference remains primitive, since its pushforward to $R$
already was primitive, and its cokernel likewise increases
by one. This proves the groups of $Q^+$. Interchanging the
sections changes the extra coordinate by a sign and does
not change any group. The dual complex has
one vertex and one self-edge, with no higher cells, and
Van Kampen again gives $\mathbb Z$.

\item In the modification $h$,
$C_{\mathrm{reg}}\simeq\mathbb C^*$ contributes zero and
the exceptional chain above the node increases Euler number
by $3-1=2$. Thus $\chi(R^\pm)=26$ and
$\chi(Q^\pm)=26-12=14$.
\end{enumerate}
\end{proof}
